This section specifies the input, processing and output path for each
experience. Simulated checkpoints use the Arduino UNO available in Tinkercad;
physical checkpoints use the MEGA2560. The diagrams cited below are included
with the corresponding development evidence. Pin numbers not fixed by the
guide must be checked against the actual sketch before they are reported as
implemented wiring.

\subsection{Sensor Systems}

\subsubsection{System 1: Analog Reading and LED Bar}
The potentiometer is used as a voltage divider between 5 V and GND, with its
wiper connected to an analog input. The program converts the 0--1023 ADC
reading into a percentage and then into 0--10 illuminated LEDs. Each LED is
placed in series with a 220- or 330-$\Omega$ resistor. The serial monitor
prints both the raw and percentage values. Tests at the minimum, midpoint
and maximum input are needed to check the scaling and the LED count. The
control flow and signal blocks are shown in Figs.~\ref{fig:2} and
\ref{fig:3}.

\subsubsection{System 2: Light and Supply Monitoring}
Four LDRs, each paired with a 10-k$\Omega$ fixed resistor, provide separate
analog light readings. The mean uses only channels that do not remain at 0
or 1023; such a channel is excluded and triggers a double buzzer warning.
When no channel remains valid, the algorithm reports a sensor fault instead
of dividing by zero. The red LED represents the luminaires and switches on
when the valid mean falls below the chosen threshold. A switch simulates the
camera supply: a drop to GND makes the blue LED blink and produces a
continuous audible alarm. The actual LDR divider orientation and threshold
must be recorded with the observed readings. See Figs.~\ref{fig:5} and
\ref{fig:6}.

\subsubsection{System 3: Occupancy and Gas Safety}
A PIR detection increments the occupancy count once, followed by a 5-s
lockout against duplicate entries. A button decrements the count for an exit;
the count is bounded below by zero and above by the configured room limit.
The LCD displays the count and an occupancy-full indication at that limit.
The MQ2 analog channel has higher priority than occupancy: above the chosen
gas threshold, the LCD displays the evacuation warning, the buzzer sounds
and the red LED blinks until the reading returns below the threshold. The
configured count limit and gas threshold need to be stated with the evidence.
The two diagrams are Figs.~\ref{fig:8} and \ref{fig:9}.

\subsubsection{System 4: Multiparameter Station}
The physical station reads distance with an HC-SR04, temperature with an
LM35, and humidity with a DHT11, then presents the values on a 2-by-16 LCD.
The normal display alternates every 3 s between distance/LM35 temperature
and humidity/``Ambiente OK'' without repeatedly clearing the entire LCD.
Distance below 10 cm or temperature above $30\,{}^{\circ}$C temporarily
overrides the second LCD line with a blinking alert, regardless of the
current view. Tests should include both normal views and each alert trigger.
The flow and blocks are Figs.~\ref{fig:10} and \ref{fig:11}.

\subsection{Actuator Systems}

\subsubsection{System 5: DC Motor in Simulation}
This exercise is simulation-only. The potentiometer ADC reading is converted
to $V=5n/1023$ V, and the motor is switched on by a digital output only when
$3<V<5$ V. Current and speed are compared for three simulated connections:
directly from the digital pin, from that pin through a 1-k$\Omega$ series
resistor, and from the 5-V pin. The resulting numbers must be taken from the
simulation instruments and their units checked. A motor must not be attached
directly to a physical Arduino output; an inductive load requires a suitable
driver and supply.

\subsubsection{System 6: PWM LED Dimming}
The potentiometer endpoints go to 5 V and GND, and its wiper goes to A0.
PWM pin 9 drives the LED anode through a 220- or 330-$\Omega$ current-limiting
resistor; the cathode returns to GND. The 10-bit ADC value is mapped to the
0--255 input of \texttt{analogWrite()}, with checks at both endpoints and an
intermediate setting. A non-PWM pin is used only for the required comparison
of its on/off behavior with gradual PWM dimming.

\subsubsection{System 7: Relay-Controlled Motor}
The relay module's VCC and GND connect to the Arduino control supply and IN
to a digital output. The external 5-V motor supply connects through COM and
the normally open (NO) contact to the motor; the motor returns to that
supply's negative terminal. If fitted, a flyback diode is placed across the
motor with its cathode on the positive side. The active HIGH/LOW level must
be checked on the actual relay module. The control and motor-power stages
are separated in Fig.~\ref{fig:15}; the motor is not powered from an Arduino
digital output.

\subsubsection{System 8: Moisture-Controlled Irrigation}
The soil probe is read in dry and wet conditions before mapping its ADC
values to a bounded 0--100\% moisture estimate. The mapping direction must
match the probe's measured response. An external supply drives the pump
through the L9110S; that supply and the Arduino share a logic-ground
reference. From 20\% to 75\%, pump, LEDs and buzzer remain off. Below 20\%,
the blue LED marks irrigation and A-IA/A-IB are HIGH/LOW for 10 s; then both
are LOW during a 15-s IDLE interval without a new measurement. Above 75\%,
the red LED and buzzer indicate an intermittent manual-attention alarm.
Calibration values, transitions and elapsed times should be recorded. The
flow and blocks are Figs.~\ref{fig:16} and \ref{fig:17}.
